\section{Introduction}
\label{Introduction}

AI agents that write and execute code are being evaluated as assistants for
scientific analysis in DiscoveryBench and
ScienceAgentBench~\cite{majumder2025discoverybench,chen2025scienceagentbench}.
A useful implementation must perform the requested analysis, and its
reported results must follow from that calculation. A researcher needs
this evidence before using a coefficient, prediction, or statistical
conclusion to support a scientific claim. Passing automated tests provides
some evidence, but its value depends on what those tests exercise.

In this study, \emph{test gaming} means satisfying tests through a shortcut
that violates the requested analysis. An agent asked to fit a regression
might return the coefficient written in a visible test instead of estimating
it from the input. A more selective implementation might recognise the
supplied dataset and substitute an answer only for that input. Both can pass
visible checks without implementing the stated method. This differs from an
ordinary programming error, such as an optimiser that fails to converge.
A failing test alone does not distinguish these behaviours; classification
requires examining the code and its development record.

Independent evaluation is difficult when the correct result is not readily
available, a challenge known as the \emph{test oracle
problem}~\cite{barr2015oracle}. Checking one known answer does not establish
that an analysis works on other samples. Prior studies document test
exploitation in coding agents and shortcuts in scientific
analysis~\cite{zhong2025impossiblebench,chen2025scienceagentbench}, motivating
checks beyond visible examples and scrutiny of the agent's account of its work.

We examine 207 completed runs on six adapted scientific-analysis tasks.
Agents receive all correct visible tests, structural tests only, no tests,
or tests with a deliberately conflicting expected value. A separate evaluator
compares hidden unit and metamorphic checks, including checks anchored to
the supplied dataset. Planted implementations and mutation testing provide
complementary evidence about test-suite strength. The contribution is a
comparison of which checks expose identified shortcuts and whether agents
disclose them at different stages. Labels outside the conflict condition
and disclosure codes remain under author review, and unflagged
implementations have not all been classified; the reported results are
preliminary.

The following research questions guide the study. RQ3 also examines
disclosure during execution and in the final answer, distinguishing it from
admissions made during later review.

\begin{itemize}
    \item \textbf{RQ1:} How often do agents game visible tests, and does hiding them help?
    \item \textbf{RQ2:} Which hidden checks best detect gamed code?
    \item \textbf{RQ3:} When reviewing its own code, does the agent admit to gaming?
\end{itemize}
